% ASTRO technical report following the compact structure of Liu et al. (2025),
% "Muon is Scalable for LLM Training". This file intentionally reports only
% experiments whose paper-facing cells are complete as of 2026-09-13.
\documentclass[10pt,twocolumn]{article}

\usepackage[margin=0.70in]{geometry}
\usepackage{amsmath,amssymb,mathtools}
\usepackage{booktabs}
\usepackage{tabularx}
\usepackage{graphicx}
\usepackage{xcolor}
\usepackage{microtype}
\usepackage{enumitem}
\usepackage[numbers,sort&compress]{natbib}
\usepackage[hidelinks]{hyperref}
\usepackage{url}
\usepackage{caption}
\usepackage{array}
\usepackage{xspace}
\usepackage{algorithm}
\usepackage{algpseudocode}
\usepackage{titlesec}

\definecolor{astropurple}{HTML}{6F3C91}
\definecolor{astrolight}{HTML}{F6F1F9}
\definecolor{commentblue}{HTML}{274CDA}
\definecolor{muted}{HTML}{666666}

\setlength{\columnsep}{0.235in}
\setlength{\parindent}{1em}
\setlength{\parskip}{0pt}
\captionsetup{font=small,labelfont=bf,skip=4pt}
\setlist{nosep,leftmargin=*}
\emergencystretch=1em
\titleformat{\section}{\bfseries\normalsize}{\thesection}{0.65em}{}
\titleformat{\subsection}{\bfseries\small}{\thesubsection}{0.55em}{}
\titlespacing*{\section}{0pt}{1.05em}{0.40em}
\titlespacing*{\subsection}{0pt}{0.85em}{0.28em}
\algrenewcommand\algorithmiccomment[1]{\hfill{\color{commentblue}\(\triangleright\) #1}}

\newcommand{\ASTRO}{\textsc{Astro}\xspace}
\newcommand{\ASTROv}{\textsc{Astro-v2}\xspace}
\newcommand{\Muon}{\textsc{Muon}\xspace}
\newcommand{\NorMuon}{\textsc{NorMuon}\xspace}
\newcommand{\AdaMuon}{\textsc{AdaMuon}\xspace}
\newcommand{\norm}[1]{\lVert #1 \rVert}
\newcommand{\good}[1]{\textcolor{astropurple}{\textbf{#1}}}
\newcommand{\finding}[1]{%
  \medskip\noindent\fcolorbox{astropurple!38}{astrolight}{%
  \parbox{\dimexpr\columnwidth-2\fboxsep-2\fboxrule}{%
  \textbf{Finding.} #1}}\medskip}
\newcommand{\paperfig}[2][\linewidth]{%
  \IfFileExists{../../artifacts/figures/#2.pdf}%
    {\includegraphics[width=#1]{../../artifacts/figures/#2.pdf}}%
    {\fbox{\parbox[c][1.2in][c]{0.94\linewidth}{\centering\small Missing generated figure: \texttt{#2}}}}}
% Keep the first-page figure behind a no-argument macro. TeX scans the optional
% argument to \twocolumn before expanding macros, so nested [...] from
% \includegraphics would otherwise terminate that argument early.
\newcommand{\herofigure}{\includegraphics[width=0.95\textwidth]{../../artifacts/figures/fig_optimizer_journey.pdf}}

\begin{document}

% ---------------------------------------------------------------------------
% First page: compact technical-report header, abstract, and headline figure.
% ---------------------------------------------------------------------------
\twocolumn[
\begin{@twocolumnfalse}
\begin{center}
{\LARGE\bfseries ASTRO: Matrix-Structured Optimization with\\
Post-Spectral Adaptive Redistribution\par}
\vspace{0.42em}
{\small\bfseries TECHNICAL REPORT\par}
\vspace{0.58em}
{\normalsize Pop Alexandru\par}
{\small Independent Researcher \quad \url{https://github.com/pop123-ux/ASTRO}\par}
{\small September 2026\par}
\end{center}
\vspace{0.35em}

\begin{minipage}{0.96\textwidth}
\small
\textbf{ABSTRACT}

\ASTRO is a matrix-structured optimizer research project built around a simple
question: after a Muon-style spectral transform has normalized matrix geometry,
can adaptive information be reintroduced without destroying the scale and
structure that made the spectral step useful? The current design first forms a
Muon-like polar direction, then maintains a row-wise second moment on that
\emph{post-polar} direction, redistributes update energy across output neurons,
and restores the original Frobenius norm. Fused attention projections may be
processed as separate Q, K, and V operators, while embeddings, tied heads,
norms, and biases use a separate Adam-like path.

Among the paper-grade ASTRO recipes measured so far, \ASTROv is the strongest.
It uses the same \texttt{Astro} implementation as the base recipe but changes
$\beta_1$ from $0.95$ to $0.9$ and disables cautious weight decay. At GPT-2
124M and 900 training steps, a completed five-configuration shared grid gives a
mean paired validation-loss difference of \good{-0.0816} versus Muon, with
\ASTROv better in \good{3/5} settings. A separate three-configuration mechanism
study gives mean paired deltas of \good{-0.2208} for \ASTROv, $-0.2137$ for the
beta-aligned ASTRO-MB flavor, and $-0.2047$ when the post-polar variance power is
removed. The same experiment measures approximately \textbf{6.4\%} additional
wall-clock per 900-step \ASTROv run. We therefore report a step-matched loss
advantage in these completed experiments, not a compute-efficiency claim.
\end{minipage}

\vspace{0.72em}
\begin{center}
\herofigure
\captionof{figure}{\textbf{Current ASTRO evidence against modern Muon-family
baselines.} Left: a pointwise shared-configuration comparison at 124M/900.
Right: paired deltas versus Muon across the completed targeted mechanism study.
Configurations are paired hyperparameter settings, not independent seeds.}
\label{fig:headline}
\end{center}
\vspace{0.45em}
\end{@twocolumnfalse}
]

\section{Motivation}

Muon demonstrated that hidden neural-network weights can be optimized as
matrices rather than as unrelated scalar coordinates. Its central operation
approximately orthogonalizes momentum with a short Newton--Schulz iteration,
thereby controlling the singular structure of the update
\citep{jordan2024muon,liu2025scalable}. Later methods such as \NorMuon and
\AdaMuon add adaptivity after orthogonalization \citep{li2025normuon,si2025adamuon}.

ASTRO was motivated by three practical observations from this line of work.
First, a spectral transform answers a geometric question---which matrix
directions should dominate---but does not by itself decide how update energy
should be distributed across different output neurons. Second, fused operators
such as GPT-2's combined QKV projection contain semantically different linear
maps inside one tensor; treating the whole tensor as one operator can allocate
update mass very unevenly between Q, K, and V. Third, a language model is not
made only of hidden matrices: embeddings, tied output heads, biases, and
normalization parameters still require an elementwise optimizer path, and
changes to that path can materially affect a supposedly ``spectral'' comparison.

The main ASTRO idea is therefore to separate \emph{geometry} from
\emph{redistribution}. ASTRO first performs the matrix-geometric cleanup, then
uses an adaptive statistic on the cleaned direction, and finally restores the
original global norm so that adaptivity changes primarily \emph{where} the step
is spent rather than silently changing the total step magnitude. Structural
routing and optional QKV splitting are part of the same principle: the optimizer
should respect what a tensor represents, not only its storage shape.

The project also adopted a methodological motivation. Several early ASTRO
results changed when learning-rate calibration, scalar-path behavior, or weight
decay changed. For that reason this report includes only experiments whose
planned paper-facing cells are complete. Incomplete seed replication, 2700-step
runs, and scale experiments are deliberately not used as findings in this
snapshot.

\section{Method}

\subsection{Muon background}

For a matrix gradient $G_t$, Muon maintains momentum and applies an approximate
polar transform. In simplified form,
\begin{align}
M_t &= \mu M_{t-1} + (1-\mu)G_t, \\
Q_t &= \operatorname{PolarApprox}(M_t), \\
W_t &= W_{t-1} - \eta_t Q_t.
\end{align}
The practical implementation uses a small number of quintic Newton--Schulz
iterations rather than an exact SVD. Large-scale Muon work further emphasizes
weight decay and per-matrix update scaling as necessary parts of a practical
LLM optimizer \citep{liu2025scalable}.

\subsection{ASTRO's spectral path}

ASTRO keeps the matrix-first principle but applies adaptive redistribution after
the polar step. Momentum is
\begin{equation}
M_t = \beta_1 M_{t-1} + (1-\beta_1)G_t.
\end{equation}
A Nesterov-style lookahead is sent through the polar iteration, optionally
block-by-block for fused QKV projections, producing $Q_t$. ASTRO then maintains
a row-wise second moment
\begin{equation}
v_t = \beta_2 v_{t-1} + (1-\beta_2)\operatorname{rowmean}(Q_t^2),
\end{equation}
with bias-corrected $\widehat v_t$. The adaptive direction is
\begin{equation}
R_t = Q_t \oslash \left(\widehat v_t^{\gamma/2}+\epsilon\right).
\end{equation}
ASTRO restores the Frobenius norm,
\begin{equation}
D_t = R_t\frac{\norm{Q_t}_F}{\norm{R_t}_F},
\end{equation}
then applies Muon-style aspect-ratio scaling per block. This norm restoration is
important conceptually: the second moment is intended to redistribute the step
across rows rather than become an uncontrolled extra learning-rate multiplier.

\subsection{Parameter routing}

The reference implementation subclasses \texttt{torch.optim.Optimizer}, but the
update rule itself is implemented explicitly in \texttt{scripts/astro\_lab.py}.
Hidden two-dimensional operators use the spectral path. Embeddings, tied output
heads, biases, normalization parameters, and other non-operators use an
Adam-like scalar path. GPT-2 \texttt{Conv1D} weights are viewed in operator
orientation before row statistics and QKV splitting are applied.

\begin{algorithm}[t]
\caption{ASTRO-v2 update}
\label{alg:astro}
\footnotesize
\begin{algorithmic}[1]
\Require parameter $W$, gradient $G$, momentum $M$, row variance $v$, step $t$, $\beta_1,\beta_2,\eta,\lambda,\gamma$
\If{$W$ is not a hidden matrix operator}
  \State $W \gets \operatorname{AdamLikeStep}(W,G,\eta\cdot s,\beta_1,\beta_2,\lambda)$
  \State \Return $W$ \Comment{$s$ is the scalar-path LR multiplier}
\EndIf
\State $G \gets \operatorname{OperatorView}(G)$ \Comment{transpose GPT-2 Conv1D storage if needed}
\State $M \gets \beta_1 M + (1-\beta_1)G$
\State $H \gets (1-\beta_1)G + \beta_1 M$ \Comment{Nesterov-style lookahead}
\State $\{H_j\} \gets \operatorname{SplitQKV}(H)$ \Comment{when the projection is fused}
\For{each block $H_j$}
  \State $Q_j \gets \operatorname{NewtonSchulz}(H_j)$
\EndFor
\State $Q \gets \operatorname{Concat}(Q_1,\ldots,Q_k)$
\State $v \gets \beta_2 v + (1-\beta_2)\operatorname{rowmean}(Q^2)$
\State $\widehat v \gets v/(1-\beta_2^t)$
\State $R \gets Q \oslash (\widehat v^{\gamma/2}+\epsilon)$
\State $D \gets R\,\norm{Q}_F/(\norm{R}_F+\epsilon)$ \Comment{preserve global update norm}
\State $D \gets \operatorname{AspectScale}(D)$ \Comment{Muon-style per-block scale}
\State $W \gets (1-\eta\lambda)W - \eta D$ \Comment{plain decoupled WD in ASTRO-v2}
\State \Return $W$
\end{algorithmic}
\end{algorithm}

\subsection{ASTRO flavors}

ASTRO is one optimizer implementation with controlled recipe variants rather
than a collection of unrelated optimizers. Table~\ref{tab:flavors} lists the
flavors currently defined in the benchmark harness. The base ASTRO defaults are
$\beta=(0.95,0.95)$, post-polar row adaptation with $\gamma=1$, QKV splitting,
Muon-style aspect scaling, and cautious weight decay.

\begin{table*}[t]
\centering
\caption{ASTRO flavors defined in the current research harness. A flavor changes
one or a small number of switches on the same \texttt{Astro} implementation.
``Measured'' means the flavor appears in the current paper-grade comparisons;
other entries are defined ablations and are not ranked here.}
\label{tab:flavors}
\scriptsize
\begin{tabularx}{\textwidth}{@{}l l X l@{}}
\toprule
Family & Flavor & Difference from base ASTRO & Paper status \\
\midrule
Reference & ASTRO & Base recipe: $\beta=(.95,.95)$, $\gamma=1$, QKV split, cautious WD & historical control \\
Calibration & ASTRO-Pinned & Set \texttt{post\_normalize=True} after each polar block & defined ablation \\
Scaling & ASTRO-Trust & Replace Muon shape scaling with layer-norm trust scaling & defined ablation \\
Masking & ASTRO-Cautious & Apply a coordinate-wise cautious update mask & defined ablation \\
Polar & ASTRO-Converging & Use the converging per-step polar polynomial schedule & defined ablation \\
Adaptation & ASTRO-$\gamma0$ & Set \texttt{variance\_power=0} & defined ablation \\
Adaptation & ASTRO-$\gamma25$ & Set \texttt{variance\_power=0.25} & defined ablation \\
Adaptation & ASTRO-$\gamma50$ & Set \texttt{variance\_power=0.5} & defined ablation \\
Geometry & ASTRO-Eq & Equalize lookahead row norms before the polar step & defined ablation \\
Decay & ASTRO-Plain-WD & Disable cautious weight decay & defined ablation \\
Decay & ASTRO-WD-Rescaled & Rescale the surviving cautious-decay coordinates & defined ablation \\
Momentum & ASTRO-MB & Set $\beta=(.9,.95)$; keep cautious WD & measured \\
Main & \textbf{ASTRO-v2} & Set $\beta=(.9,.95)$ and disable cautious WD & \textbf{best measured flavor} \\
Attribution & ASTRO-v2-$\gamma0$ & ASTRO-v2 with \texttt{variance\_power=0} & measured \\
Structure & ASTRO-NoSplit & Disable QKV block splitting & defined ablation \\
Structure & ASTRO-Split100 & Split QKV only through step 100 & defined ablation \\
Structure & ASTRO-Split300 & Split QKV only through step 300 & defined ablation \\
\bottomrule
\end{tabularx}
\end{table*}

\subsection{Why ASTRO-v2 is the main flavor}

\ASTROv is not a separate class. It changes only two base-recipe switches:
\begin{equation}
(\beta_1,\beta_2)=(0.9,0.95), \qquad \text{cautious WD}=\text{off}.
\end{equation}
The first change aligns the shared beta pair with the behavior used by Muon's
auxiliary Adam-like path instead of forcing $\beta_1=0.95$ onto both spectral
and scalar groups. The second removes a masked decay rule that changes the
effective amount of regularization. ASTRO-MB isolates the beta change, while
ASTRO-v2-$\gamma0$ asks whether the post-polar variance term is carrying the
result.

Among flavors with directly comparable paper-grade measurements, ASTRO-v2 is
currently the best: it has the lowest loss in the one shared modern-baseline
comparison and the most negative mean paired delta in the targeted mechanism
study. This does \emph{not} rank unmeasured flavors.

\section{Experiments}\par\vspace{0.45em}

\subsection{Experimental setup}

All experiments reported in this section are already complete. They use
GPT-2-style causal language models trained from scratch on FineWeb-Edu with
GPT-2 BPE. The principal model has approximately 124M parameters, 12 layers, 12
attention heads, width 768, and sequence length 512. Validation is pinned to the
same first 400,000 tokens. The principal metric is validation cross-entropy
loss; runtime experiments additionally report wall-clock seconds.

Two forms of comparison are kept distinct. A \emph{shared configuration}
means optimizers receive the same learning rate, weight decay, and scalar-path
multiplier. It probes sensitivity and robustness, not random-seed significance.
The report does not use incomplete independent-seed or scale experiments.

\subsection{Experiment 1: tuned 300-step control}

The first experiment asks whether the earlier ASTRO recipe already beats strong
baselines under an expanded tuned short-horizon protocol. It does not.

\begin{table}[t]
\centering
\caption{GPT-2 124M, 300 steps, three evaluation seeds. The ASTRO row is the
earlier recipe rather than ASTRO-v2. Lower is better.}
\label{tab:300}
\small
\begin{tabular}{lrr}
\toprule
Optimizer & Mean val. loss & $\Delta$ vs Muon \\
\midrule
NorMuon & \textbf{6.6280} & \textbf{-0.0077} \\
Muon & 6.6357 & 0.0000 \\
Earlier ASTRO & 6.6538 & +0.0180 \\
AdamW & 6.8445 & +0.2087 \\
\bottomrule
\end{tabular}
\end{table}

\finding{The earlier ASTRO recipe is worse than Muon by $+0.0180$ validation
loss in this tuned 300-step benchmark. ASTRO-v2 should therefore be treated as a
later recipe, not retroactively credited for earlier ASTRO runs.}

\subsection{Experiment 2: modern Muon-family baselines at 124M/900}

The second experiment compares Muon, NorMuon, AdaMuon, ASTRO-MB, and ASTRO-v2 at
one identical configuration:
$\eta=0.0102059788$, weight decay $0.0010597179$, and scalar multiplier $0.4369$.

\begin{table}[t]
\centering
\caption{One shared GPT-2 124M/900 configuration. This is a pointwise
comparison, not an independently tuned global ranking.}
\label{tab:shared}
\small
\begin{tabular}{lrr}
\toprule
Optimizer & Val. loss & $\Delta$ vs Muon \\
\midrule
\textbf{ASTRO-v2} & \textbf{5.9208} & \textbf{-0.0827} \\
ASTRO-MB & 5.9226 & -0.0809 \\
AdaMuon & 6.0028 & -0.0007 \\
Muon & 6.0035 & 0.0000 \\
NorMuon & 6.0155 & +0.0121 \\
\bottomrule
\end{tabular}
\end{table}

\finding{ASTRO-v2 is the lowest-loss method in this shared setting, but
ASTRO-MB is only 0.0019 worse. The near-tie immediately suggests that the full
post-polar adaptation is not the only plausible source of the gain.}

\subsection{Experiment 3: targeted ASTRO mechanism study}

The third experiment holds the scalar multiplier at $0.4369$ and evaluates
Muon, ASTRO-MB, ASTRO-v2, and ASTRO-v2-$\gamma0$ across three shared
hyperparameter settings.

\begin{table}[t]
\centering
\caption{Completed targeted 124M/900 mechanism study. Configurations are paired
settings, not independent seeds.}
\label{tab:mechanism}
\small
\begin{tabular}{lrr}
\toprule
Flavor & Mean $\Delta$ vs Muon & Settings won \\
\midrule
\textbf{ASTRO-v2} & \textbf{-0.2208} & \textbf{3/3} \\
ASTRO-MB & -0.2137 & 3/3 \\
ASTRO-v2-$\gamma0$ & -0.2047 & 3/3 \\
\bottomrule
\end{tabular}
\end{table}

The largest separation occurs at $\eta\approx0.0505$: Muon reaches 6.4527
validation loss while ASTRO-v2 remains at 5.9436. At the other two settings the
gaps are substantially smaller.

\finding{Most of ASTRO-v2's measured advantage survives when $\gamma=0$ and is
already present in ASTRO-MB. The completed evidence therefore supports a
robustness interpretation more strongly than a claim that variance-power
adaptation alone is the mechanism.}

\subsection{Experiment 4: five-setting 900-step robustness grid}

A separate completed five-configuration grid evaluates Muon and ASTRO-v2 at the
same five settings. Figure~\ref{fig:paired} shows every pair.

\begin{figure}[t]
\centering
\paperfig{fig_paired_900}
\caption{\textbf{All five shared 124M/900 configurations.} Each grey connector
joins Muon and ASTRO-v2 under identical hyperparameters. The mean paired delta
is $-0.0816$ and 3/5 settings favor ASTRO-v2.}
\label{fig:paired}
\end{figure}

The five paired deltas are approximately $-0.091$, $+0.067$, $-0.108$,
$+0.006$, and $-0.283$. The mean is $-0.0816$. The best ASTRO-v2 loss observed
inside this five-setting grid is 5.8102; the best Muon loss is 5.9286. Those two
minima come from different configurations and are not presented as a paired
sample.

\finding{ASTRO-v2 improves the mean across the five shared settings and wins
three of them, but the effect is not uniform. The shape of the result is
consistent with a wider stable hyperparameter region rather than dominance at
every point.}

\subsection{Experiment 5: training-horizon study through 900 steps}

The fifth experiment repeats the same five shared configurations at 300, 600,
and 900 steps. Only the three complete horizons are reported here.

\begin{table}[t]
\centering
\caption{Completed ASTRO-v2 horizon cells. Negative mean delta favors ASTRO-v2.}
\label{tab:horizon}
\small
\begin{tabular}{rrrr}
\toprule
Steps & Settings & Mean $\Delta$ & ASTRO-v2 wins \\
\midrule
300 & 5/5 & -0.0500 & 2/5 \\
600 & 5/5 & -0.0541 & 3/5 \\
900 & 5/5 & \textbf{-0.0816} & 3/5 \\
\bottomrule
\end{tabular}
\end{table}

\begin{figure}[t]
\centering
\paperfig{fig_horizon_v2}
\caption{Paired ASTRO-v2 minus Muon validation-loss deltas over the three
completed horizons. Small points are individual shared settings; diamonds are
cell means.}
\label{fig:horizon}
\end{figure}

\finding{The mean paired margin is negative at 300, 600, and 900 steps and does
not collapse by 900 steps. This establishes persistence through the completed
900-step horizon, but says nothing about longer unmeasured training.}

\subsection{Experiment 6: fused QKV allocation}

The sixth experiment probes how the spectral transform allocates update mass
inside fused QKV projections. The same tensor is treated either as one fused
operator or as three Q/K/V blocks before the polar step.

\begin{table}[t]
\centering
\caption{Normalized update allocation across Q/K/V. Splitting is applied before
the polar transform.}
\label{tab:qkv}
\small
\begin{tabular}{lcc}
\toprule
Checkpoint & Fused Q/K/V & Split Q/K/V \\
\midrule
GPT-2 & .244/.226/\textbf{.530} & .318/.347/.336 \\
GPT-2 Medium & .245/.266/\textbf{.489} & .308/.369/.323 \\
Pythia-410M & .218/.248/\textbf{.534} & .356/.300/.344 \\
\bottomrule
\end{tabular}
\end{table}

The corresponding gradient V/Q ratios are approximately $2.65\times$,
$3.27\times$, and $3.28\times$. Random-initialization controls also show V
skew, so training is not claimed as the sole cause.

\finding{Fused spectral treatment produces strongly non-uniform Q/K/V update
allocation in these probes, while block-wise treatment makes the allocation
substantially more balanced.}

\subsection{Experiment 7: same-step runtime cost}

The targeted mechanism study also records wall-clock time. Muon averages roughly
787 seconds per 900-step run and ASTRO-v2 roughly 838 seconds, an overhead of
approximately 6.4\%.

\begin{figure}[t]
\centering
\paperfig{fig_efficiency_v2}
\caption{Validation loss versus measured wall-clock in the completed targeted
124M/900 experiment. Thin connectors join methods evaluated at the same setting.
This is same-step timing, not a time-matched training study.}
\label{fig:runtime}
\end{figure}

\finding{ASTRO-v2's current advantage is a better-loss-at-equal-steps result.
Because each step is measurably more expensive, the completed experiments do not
establish a compute-efficiency improvement.}

\section{Discussion}

The completed evidence supports two conclusions. First, ASTRO-v2 is the best
measured ASTRO flavor in the current controlled comparisons. It improves over
Muon on average across the completed five-setting 900-step grid and remains the
strongest of ASTRO-MB, ASTRO-v2, and ASTRO-v2-$\gamma0$ in the targeted study.
Second, the mechanism is simpler than the original design story suggested.
ASTRO-MB nearly matches ASTRO-v2, and removing variance power preserves most of
the gain. The data therefore point toward beta/scalar-path alignment, decay
behavior, and robustness under aggressive settings as important parts of the
current result.

This does not make the post-polar ASTRO idea irrelevant. The optimizer is still
organized around a useful separation: spectral geometry first, adaptive
redistribution second, norm restoration afterward. What the ablation says is
that the most distinctive mathematical component should not be credited with a
performance gain merely because it is distinctive.

The QKV probe provides a separate structural result. It shows why ASTRO treats
fused attention projections as potentially composite operators: the polar
transform can be mathematically well-defined on the fused matrix while
allocating update mass very unevenly between the three semantic sub-operators.

\section{Limitations}

This snapshot is intentionally narrower than the full research program. The
main training evidence is GPT-2 124M. The completed ASTRO-v2 robustness grids
pair hyperparameter configurations rather than independent random seeds. The
900-step advantage is therefore evidence about behavior across settings, not a
statistical seed-replication claim. ASTRO-v2 also costs about 6.4\% more
wall-clock per step in the completed targeted study. Results from incomplete
longer-horizon, independent-seed, or larger-model experiments are excluded from
this report rather than extrapolated.

\section{Conclusion}

ASTRO combines Muon-style matrix geometry with structure-aware routing and
post-spectral adaptive redistribution. Its core operation is deliberately
ordered: orthogonalize first, adapt the cleaned direction second, and restore
its global norm before applying the update. ASTRO-v2 is currently the strongest
measured recipe of this family. At 124M/900 it achieves a mean paired delta of
$-0.0816$ across five shared settings and wins 3/5; in a separate targeted
three-setting study it records the strongest mean delta among the measured
ASTRO flavors while carrying about 6.4\% runtime overhead.

The clearest interpretation of the completed evidence is not universal
optimizer dominance. It is that ASTRO-v2 can remain stable in parts of the
hyperparameter region where Muon degrades, and that much of this behavior is
already present in the simpler ASTRO-MB flavor. The paper therefore treats
ASTRO-v2 as the current main recipe while preserving the full flavor family as
an attribution framework.

\paragraph{Reproducibility.}
The implementation, experiment ledger, paper-facing measurements, plotting
scripts, and manuscript source are available at
\url{https://github.com/pop123-ux/ASTRO}.

\begin{thebibliography}{99}
\bibitem[Jordan et al.(2024)]{jordan2024muon}
K.~Jordan, Y.~Jin, V.~Boza, J.~You, F.~Cesista, L.~Newhouse, and J.~Bernstein.
\newblock Muon: An optimizer for hidden layers in neural networks.
\newblock 2024. \url{https://kellerjordan.github.io/posts/muon/}.

\bibitem[Liu et al.(2025)]{liu2025scalable}
J.~Liu, J.~Su, X.~Yao, et al.
\newblock Muon is Scalable for LLM Training.
\newblock \emph{arXiv:2502.16982}, 2025.

\bibitem[Li et al.(2025)]{li2025normuon}
Z.~Li et al.
\newblock NorMuon: Making Muon more efficient and scalable.
\newblock \emph{arXiv:2510.05491}, 2025.

\bibitem[Si et al.(2025)]{si2025adamuon}
C.~Si, D.~Zhang, and W.~Shen.
\newblock AdaMuon: Adaptive Muon Optimizer.
\newblock \emph{arXiv:2507.11005}, 2025.

\bibitem[Wen et al.(2025)]{wen2025fantastic}
K.~Wen, D.~Hall, T.~Ma, and P.~Liang.
\newblock Fantastic Pretraining Optimizers and Where to Find Them.
\newblock \emph{arXiv:2509.02046}, 2025.
\end{thebibliography}

\end{document}
